\section{API Rate Limiting zum Schutz deterministischer Blind Indizes (5123034, 5123060, \\5123097)}

Mit der Einführung des Blind-Indexing-Verfahrens zur performanten Filterung verschlüsselter Attribute wurde eine deterministische Komponente in die Datenhaltung integriert. Da identische Klartextwerte (z.B. der Nachname "`Müller"') stets zu identischen Hashwerten führen, ergibt sich hieraus ein theoretischer Angriffsvektor: Erlangt ein Angreifer lesenden Zugriff auf die Datenbank oder missbraucht die API, könnte er durch automatisiertes Durchprobieren gängiger Werte (Wörterbuch-Attacke / Brute-Force) und den anschließenden Abgleich der Hashes Rückschlüsse auf die echten Daten ziehen.

Dieses Risiko betrifft grundsätzlich sämtliche Domänenobjekte, deren sensible Attribute verschlüsselt gespeichert und über Blind Indizes durchsuchbar gemacht werden. Dazu zählen in der vorliegenden Anwendung insbesondere Patienten-, Pflegekraft-, Arzt-, Diagnose-, Medikations- und Buchungsdaten. Um dieses Risiko zu minimieren, wurde für alle entsprechenden Such- und Filterendpunkte ein restriktives \textit{API Rate Limiting} auf Applikationsebene implementiert, welches die Frequenz der ausführbaren Suchanfragen begrenzt.

\subsection{Technische Implementierung}

Die Umsetzung des Ratenbegrenzungs-Mechanismus basiert auf dem etablierten Token-Bucket-Algorithmus und wurde wie folgt in die bestehende Spring-Boot-Architektur integriert:

\begin{enumerate}

\item \textbf{Zentraler Interceptor (\texttt{RateLimitingInterceptor.kt}):} Die Kernlogik verwaltet die Buckets zustandslos und thread-sicher über eine \texttt{ConcurrentHashMap<String, Bucket>}, wobei die IP-Adresse des anfragenden Clients als Schlüssel dient. Jeder Client-IP wird ein Kontingent von maximal $10$ Tokens zugewiesen, welches sich nach dem "`greedy refill"'-Prinzip alle 60 Sekunden vollständig erneuert. Ist das Token-Kontingent erschöpft, bricht die Anfrage sofort ab und liefert den HTTP-Status \texttt{429 Too Many Requests} mitsamt einer standardisierten JSON-Fehlermeldung zurück.

\item \textbf{Spezifische Routen-Registrierung (\texttt{WebConfig.kt}):} Der Interceptor wurde über die \texttt{WebMvcConfigurer}-Klasse registriert. Der Schutz wird gezielt auf alle Endpunkte angewendet, die Such- oder Filteroperationen über verschlüsselte Daten mittels Blind Indizes ermöglichen. Hierzu zählen insbesondere die Ressourcen \texttt{/api/patients}, \texttt{/api/nurses}, \texttt{/api/doctors}, \texttt{/api/diagnoses}, \texttt{/api/medications} sowie \texttt{/api/bookings}. Administrative Endpunkte oder rein schreibende Operationen (z.\,B. das Anlegen neuer Datensätze) sind hiervon nicht betroffen, da dort kein Risiko systematischer Blind-Index-Angriffe besteht.

\end{enumerate}


\subsection{Wahl der Rate-Limiting-Parameter}

Für die vorliegende Anwendung wurde ein einheitliches Rate Limit von
10 Anfragen pro Minute und Client-IP umgesetzt. Nach Ablauf einer Minute
wird das verfügbare Token-Kontingent automatisch wieder aufgefüllt.
Dieser Wert wurde bewusst niedrig gewählt, um die Funktionsweise des
Schutzmechanismus während der Entwicklung und Validierung leicht
nachvollziehbar zu machen und Wörterbuch-Angriffe auf die verwendeten
Blind Indizes wirksam einzuschränken.

In einem produktiven Krankenhausinformationssystem würden die konkreten
Grenzwerte abhängig vom erwarteten Nutzerverhalten und der
Systemauslastung festgelegt werden. Dabei wären deutlich höhere Limits
denkbar, beispielsweise mehrere hundert oder tausend Anfragen innerhalb
eines definierten Zeitraums. Solche Werte liegen in der Regel weit über
dem Anfragevolumen eines normalen Benutzers, können jedoch bei
automatisierten Angriffen erreicht werden.

Zusätzlich wäre es möglich in produktiven Umgebungen weitergehende
Maßnahmen einzusetzen, etwa längere Sperrzeiten nach Überschreiten des
Limits, gestaffelte Strafzeiten oder die temporäre Blockierung
auffälliger Clients. Dadurch steigt der Aufwand für automatisierte
Wörterbuch- und Brute-Force-Angriffe erheblich, während legitime Nutzer
im Normalbetrieb nicht beeinträchtigt werden.

\subsection{Validierung und Verifikation}
Die Funktion der Ratenbegrenzung wurde mittels automatisierter HTTP-Anfragen im Terminal verifiziert. Dabei simuliert eine Schleife zwölf sequentielle Aufrufe des geschützten Such-Endpunkts in direkter Abfolge:

\begin{verbatim}
for i in $(seq 1 12); do
  echo -n "Request $i: "
  curl -s -o /dev/null -w "%{http_code}" "http://localhost:8080/api/patients"
  echo
done
\end{verbatim}

für MacOS bzw. unter Windows:

\begin{verbatim}
for /L %i in (1,1,12) do (
    echo Request %i:
    curl -s -o NUL -w "%{http_code}" http://localhost:8080/api/patients
    echo.
)
\end{verbatim}

Das Testergebnis bestätigt die korrekte Funktionsweise der Middleware: Die ersten zehn Anfragen werden erfolgreich verarbeitet (HTTP 200), während die Anfragen 11 und 12 aufgrund des erschöpften Token-Buckets blockiert werden (HTTP 429).


